\chapter{Reference: colour and style modifiers}
\label{ch:refman-style}

\begin{aside}
The full palette of inline style. Apply with the pipe
operator; chain as many as you like; the order is mostly
irrelevant.
\end{aside}

\section{Foreground / background}

\begin{itemize}[leftmargin=*]
  \item \code{Fg<R, G, B>}: 24-bit foreground.
  \item \code{Bg<R, G, B>}: 24-bit background.
  \item \code{Fg256<N>}, \code{Bg256<N>}: 256-colour
    palette index.
  \item \code{Fg16<Red>}, \code{Bg16<Blue>}: legacy 16-colour
    names.
\end{itemize}

\section{Weight and decoration}

\begin{itemize}[leftmargin=*]
  \item \code{Bold}, \code{Dim}, \code{Italic},
    \code{Underline}.
  \item \code{Blink}, \code{Reverse}, \code{Strike},
    \code{Hidden}.
\end{itemize}

\section{Working example}

\begin{lstlisting}[language=C++]
#include <maya/maya.hpp>

using namespace maya;
using namespace maya::dsl;

struct StyleDemo {
    struct Model {}; struct Quit {};
    using Msg = std::variant<Quit>;

    static Model init() { return {}; }
    static auto update(Model m, Msg)
        -> std::pair<Model, Cmd<Msg>> {
        return {m, Cmd<Msg>::quit()};
    }

    static Element view(const Model&) {
        return v(
            t<"bold">              | Bold,
            t<"dim">               | Dim,
            t<"italic">            | Italic,
            t<"underline">         | Underline,
            t<"strike">            | Strike,
            t<"reverse">           | Reverse,
            blank_,
            t<"red on black">      | Fg<255, 80, 80>,
            t<"green on dark">     | Fg<100, 220, 120>,
            t<"blue on dark">      | Fg<100, 180, 255>,
            t<"warning">           | Fg<20, 20, 20> | Bg<255, 220, 80>,
            blank_,
            t<"bold + italic + cyan"> | Bold | Italic | Fg<100, 220, 220>,
            blank_,
            t<"q to quit">         | Dim
        ) | pad<1> | border_<Round>;
    }

    static auto subscribe(const Model&) -> Sub<Msg> {
        return key_map<Msg>({ {'q', Quit{}} });
    }
};

int main() { run<StyleDemo>({.title = "style"}); }
\end{lstlisting}

\section{Notes}

\begin{itemize}[leftmargin=*]
  \item Modifiers compose left-to-right; later ones win on
    conflicts (e.g. two foregrounds).
  \item True colour is the default; the runtime
    automatically degrades to 256-colour where the terminal
    cannot do 24-bit.
  \item \code{Blink} is honoured rarely; treat as a hint.
\end{itemize}

\section{Recap}

\begin{recap}
\begin{itemize}[leftmargin=*]
  \item Fg / Bg in 24-bit, 256, or 16.
  \item Bold, Dim, Italic, Underline, Strike, Reverse.
  \item Chain with pipes.
\end{itemize}
\end{recap}

\section*{Workshop}

\begin{enumerate}[leftmargin=*]
  \item Build a status bar with three coloured pills.
  \item Use \code{Reverse} to highlight a selected row.
  \item Make a "danger" badge: bold, white-on-red.
\end{enumerate}
